% models/mathematical/zeitfeld_formalismus.tex
% Formaler Grundaufbau des integrierten Zeitfeldes

\section{Formaler Zeitfeld-Formalismus}

\subsection{Definition des Zeitfeld-Tensors}

Das integrierte Zeitfeld wird durch einen Tensor zweiter Stufe definiert:
\[
\mathcal{T}_{\mu\nu}(x) = 
\mathcal{T}^{(p)}_{\mu\nu}(x) +
\mathcal{T}^{(s)}_{\mu\nu}(x) +
\mathcal{T}^{(g)}_{\mu\nu}(x) +
\mathcal{T}^{(r)}_{\mu\nu}(x) +
\mathcal{T}^{(q)}_{\mu\nu}(x)
\]
mit den Indizes:
\[
p = \text{planetar}, \quad
s = \text{solar}, \quad
g = \text{galaktisch}, \quad
r = \text{relativistisch}, \quad
q = \text{quantenphysikalisch}.
\]

Jede Schicht besitzt eine eigene Metrikstruktur und eigene Dynamik.

\subsection{Zeitgradienten}

Die lokalen Zeitgradienten ergeben sich aus:
\[
\nabla_\alpha T = \partial_\alpha T(x)
\]
wobei $T(x)$ die effektive Zeitrate am Ort $x$ beschreibt.

\subsection{Interferenzstruktur}

Die Interferenz zweier Zeitfelder $T_a$ und $T_b$ ergibt:
\[
T_{\text{int}} = T_a + T_b + 2\sqrt{T_a T_b}\cos(\Delta\phi)
\]
mit der Phasendifferenz $\Delta\phi$.

\subsection{Gesamtzeitrate}

Die beobachtete Zeitrate ergibt sich aus:
\[
\frac{d\tau}{dt} = f\big(
T^{(p)}, T^{(s)}, T^{(g)}, T^{(r)}, T^{(q)}
\big)
\]
wobei $f$ eine nichtlineare, schichtübergreifende Kopplungsfunktion ist.

\subsection{Kopplungsoperator}

Die Kopplungen zwischen den Schichten werden durch einen Operator $\mathcal{K}$ beschrieben:
\[
\mathcal{K}\big(\mathcal{T}^{(i)}, \mathcal{T}^{(j)}\big)
= \lambda_{ij} \, \mathcal{T}^{(i)}_{\mu\nu} \mathcal{T}^{(j)\,\mu\nu}
\]
mit Kopplungsparametern $\lambda_{ij}$.

\subsection{Fundamentale Schicht}

Die quantenphysikalische Schicht definiert die fundamentale Zeitrichtung:
\[
\mathcal{T}^{(q)}_{00} = \langle \Psi | \hat{H} | \Psi \rangle
\]
und koppelt über Entanglement-Korrelationen an die höheren Schichten.

\subsection{Relativistische Schicht}

Die relativistische Zeitrate ergibt sich aus:
\[
d\tau = \sqrt{g_{00}} \, dt
\]
wobei $g_{00}$ durch lokale Energie- und Massenverteilung bestimmt wird.

\subsection{Galaktische Schicht}

Die galaktische Komponente ergibt sich aus:
\[
\mathcal{T}^{(g)}_{00} = \Phi_{\text{DM}}(x) + \Phi_{\text{gal}}(x)
\]
mit Dunkle-Materie-Potential $\Phi_{\text{DM}}$ und galaktischem Potential $\Phi_{\text{gal}}$.

\subsection{Solare Schicht}

Die solare Modulation wird modelliert durch:
\[
\mathcal{T}^{(s)}_{0i} = A_{\odot}(t) \, B_{\odot,i}(t)
\]
mit solarem Aktivitätsindex $A_{\odot}$ und Magnetfeldkomponenten $B_{\odot,i}$.

\subsection{Planetare Schicht}

Die planetare Schicht basiert auf:
\[
\mathcal{T}^{(p)}_{ij} = \sum_k R_k \cos(\omega_k t + \phi_k)
\]
mit planetaren Frequenzen $\omega_k$ und Phasen $\phi_k$.

\subsection{Gesamtmodell}

Das vollständige Zeitfeld ergibt sich aus:
\[
\boxed{
\mathcal{T}_{\mu\nu}(x) =
\sum_{i} \mathcal{T}^{(i)}_{\mu\nu}(x)
+
\sum_{i \neq j} \mathcal{K}\big(\mathcal{T}^{(i)}, \mathcal{T}^{(j)}\big)
}
\]

Dies bildet die mathematische Grundlage des integrierten Zeitfeldmodells.
